\chapter*{Anexos}
\addcontentsline{toc}{chapter}{Anexos}

\appendix

\section{Funciones hiperbólicas}

Existe un tipo particular de funciones que poseen algunas similitudes con las funciones trigonométricas,
las funciones hiperbólicas. Es posible que su nombre corresponda con esa semejanza o al siguiente hecho,
si $y$ corresponde al \emph{seno hiperbólico} y $x$ al \emph{coseno hiperbólico},
entonces se satisface la ecuación de la hipérbola $x^{2}-y^{2}=1$.

Entre los aspectos comparables están las principales identidades,
sus derivadas, algunos límites relevantes o en algunos de sus valores como $\sinh 0=0$, $\cosh 0=1$,
$\tanh 0=0$; también hay diferencias, pues las funciones hiperbólicas no son periódicas.
\begin{defi}
  Las siguientes expresiones definen las funciones hiperbólicas.
  \begin{eqnarray}
    \sinh x &=& \frac{e^{x}-e^{-x}}{2}=\dfrac{1}{\csch x} \label{hiper1}  \\
    \cosh x &=& \frac{e^{x}+e^{-x}}{2}=\dfrac{1}{\sech x} \label{hiper2}  \\
    \tanh x &=& \dfrac{\sinh x}{\cosh x}=\dfrac{1}{\coth x} \label{hiper3}
  \end{eqnarray}
\end{defi}

Se muestran los gráficos de algunas de ellas.

\begin{figure}[H]
  \centering
  \caption{Cono truncado oblicuo determinado por ${\bf r}(u,t)$}
  \includegraphics[width=0.4\textwidth]{Fig/190.pdf} \hspace{5mm}
  \includegraphics[width=0.4\textwidth]{Fig/191.pdf}

  \includegraphics[width=0.4\textwidth]{Fig/192.pdf}
  \fuentepropia
\end{figure}

\begin{ejer} Usando la definición de funciones hiperbólicas
  \begin{enumerate}
    \item \textit{Demostrar las siguientes identidades}
      \begin{itemize}
        \item   $\cosh x+\sinh x=e^{x}$
        \item  $\cosh x-\sinh x=e^{-x}$
        \item  $\sinh (-x)=-\sinh x$
        \item   $\cosh (-x)=\cosh x$
        \item   $\tanh (-x)=-\tanh x$
        \item  $\cosh ^{2}x-\sinh ^{2}x=1$
        \item  $1-\tanh ^{2}x=\sech^{2}x$
        \item  $1-\coth ^{2}x=-\csch^{2}x$

        \item   $\dfrac{1+\tanh x}{1-\tanh x}=e^{2x}$
        \item  $\cosh ^{2}x=\dfrac{\cosh 2x+1}{2}$
        \item  $\sinh ^{2}x=\dfrac{\cosh 2x-1}{2}$
        \item  $\dfrac{\coth x-1}{\coth x+1}=e^{-2x}$
        \item   $(\cosh x+\sinh x)^{n}=\cosh nx+\sinh nx$
        \item   $\sinh (x+h)=\sinh x\cosh h+\cosh x\sinh h$
        \item   $\cosh (x+h)=\cosh x\cosh h+\sinh x\sinh h$
        \item $\sinh (x-y)= \sinh x\cosh y-\cosh x\sinh y$
        \item $\cosh (x-y)= \cosh x\cosh y-\sinh x\sinh y$
        \item   $\tanh (x+h)=\dfrac{\tanh x+\tanh y}{1+\tanh x\tanh y}$
      \end{itemize}
    \item Pruebe que

      \begin{itemize}
        \item $\lim\limits_{x\rightarrow 0}\frac{\sinh x}{x}=1$,
        \item  $\lim\limits_{x\rightarrow 0}\frac{1-\cosh x}{x}=0$,
        \item $\lim\limits_{x\rightarrow \infty }\tanh x= 1$,
        \item  $\lim\limits_{x\rightarrow -\infty }\tanh x= -1$,
        \item $\lim\limits_{x\rightarrow \infty }\coth x= 1$,
        \item  $\lim\limits_{x\rightarrow -\infty }\coth x= -1$.
      \end{itemize}

    \item Obtenga el gráfico de las funciones $y=\coth x$, $y=\sech x$, $y=\csch x$ y de las funciones
      hiperbólicas inversas.

    \item Hallar las derivadas de las funciones hiperbólicas.
  \end{enumerate}
\end{ejer}

\begin{tcolorbox}
  A partir de la definición se establece lo siguiente
  \begin{eqnarray*}
    x &=& \sinh y=\frac{1}{2}(e^{y}+e^{-y}) \\
    &=& \frac{1}{2}\bigg (e^{y}+\dfrac{1}{e^{y}}\bigg ) =  \dfrac{e^{2y}+1}{e^{y}}
  \end{eqnarray*}
  Luego,
  \begin{eqnarray*} y &=& \arcsinh x \Longleftrightarrow e^{2y}-2x e^{y}+1=0 \\ e^{y} &=& \frac{2x\pm\sqrt[2]{4x^{2}+4}}{2}
  \end{eqnarray*}
  Por lo tanto,
  \begin{equation}\label{hiper4}
    y=\arcsinh x= \ln \left| x\pm \sqrt[2]{x^{2}+1}\right|.
  \end{equation}
\end{tcolorbox}

\begin{ejer} Resolver los siguientes problemas
  \begin{enumerate}
    \item \textit{Expresar en función de $x$, de manera semejante a la ecuación \ref{hiper4}
      las siguientes funciones:}
      \begin{enumerate}
        \item $y=\arccosh x$,
        \item $y=\arctanh x$,
        \item $y=\arccoth x$,
        \item $y=\arcsech x$,
        \item $y={\rm arccsch} x$.
      \end{enumerate}

    \item Muestre que las derivadas de las funciones hiperbólicas inversas son las siguientes, compruebe las siguientes igualdades
      \begin{itemize}
        \item $\frac{d}{dx}\left({\rm arcsinh}(x) \right) = \dfrac{1}{\sqrt{x^{2}+1}}$
        \item $\frac{d}{dx}\left({\rm arccosh}(x)\right) =\dfrac{1}{\sqrt{x^{2}-1}}$
        \item $\frac{d}{dx}\left({\rm arctanh}(x)\right) =\dfrac{1}{1-x^{2}}$
        \item $\frac{d}{dx}\left({\rm arccoth}(x)\right) =\dfrac{1}{1-x^{2}}$
        \item $\frac{d}{dx}\left({\rm arcsech}(x)\right) =\dfrac{-1}{x\sqrt{1-x^{2}}}$
        \item $\frac{d}{dx}\left({\rm arccsch}(x)\right) =\dfrac{-1}{x\sqrt{1+x^{2}}}$
      \end{itemize}
  \end{enumerate}
\end{ejer}

\section{Sumatorias}

Presentamos algunas fórmulas que se requieren para los elementos que se trabajarán de este apartado. La primera expresión corresponde a sumar $n$ veces una constante, esto es:
\begin{equation} \label{sumas}
  \sum_{k=1}^{n}c= c+c+c+\cdots +c= cn
\end{equation}
Otra expresión muy usual, es la suma de los primeros $n$ números naturales. Para esta suma, es claro que
\begin{eqnarray*}
  S &=& 1+2+3+\cdots +\left( n-2\right) +\left( n-1\right) +n \\
  S &=& n+\left( n-1\right) +\left( n-2\right) +\cdots +3+2+1.
\end{eqnarray*}
Al sumar estas dos expresiones haciendo pares de sumandos se obtiene
\begin{eqnarray*}
  2S &=& ( n+1) +(n+1)+(n+1) +\cdots +(n+1)\\
  S&=&\frac{n(n+1)}{2}
\end{eqnarray*} En resumen,
\begin{equation} \label{sumas1}
  \ds\sum_{k=1}^{n}k= \frac{1}{2}n(n+1) .
\end{equation}
La siguiente suma corresponde a un tipo llamado \emph{sumas telescópicas}. \index{Sumas telescópicas}
\[\sum_{k=1}^{n}\left( \left( k+1\right) ^{3}-k^{3}\right) =\sum_{k=1}^{n}\left( 3k^{2}+3k+1\right)\!,\]
que se obtiene al hacer la expansión y simplificar términos,
\begin{multline*}
  \left(2^{3}-1^{3}\right) +\left( 3^{3}-2^{3}\right) +\cdots +\left( n^{3}-\left(
  n-1\right) ^{3}\right) +\left( (n+1) ^{3}-n^{3}\right)\\ = (n+1) ^{3}-1;
\end{multline*}
entonces,
\begin{eqnarray*}
  \sum_{k=1}^{n}\left( 3k^{2}+3k+1\right) &=& 3\sum_{k=1}^{n}k^{2}+3\sum_{k=1}^{n}k+\sum_{k=1}^{n}1 \\
  &=& 3\sum_{k=1}^{n}k^{2}+\frac{3n(n+1) }{2}+n = (n+1)^{3}-1.
\end{eqnarray*}
De donde se desprende la expresión
\begin{equation} \label{sumas2}
  \sum_{k=1}^{n}k^{2}= \frac{1}{6}n(n+1) \left(2n+1\right)
\end{equation}
\begin{tcolorbox}
  De manera general, con la suma telescópica
  \[\sum\limits _{k=1}^{n}\left( \left(k+1\right) ^{r}-k^{r}\right) = (n+1) ^{r}-1,\]
  es posible obtener la fórmula para sumar las potencias $r-$ésimas de los primeros $n$ números naturales,
  para todo $r\in \mathbb{N}$.
\end{tcolorbox}
En particular,
\begin{eqnarray}
  \sum_{k=1}^{n}k^{3} &=& \frac{n^{2}\left(n+1\right)^{2}}{4}  \label{sumas3} \\
  \sum_{k=1}^{n}k^{4} &=& \frac{n\left( 2n+1\right) \left(n+1\right) \left( 3n^{2}+3n-1\right)}{30} \label{sumas4} \\
  \sum_{k=1}^{n}k^{5} &=& \frac{n^{2}(n+1)^2(2n^{2}+2n-1)}{12}   \label{sumas5}\\
  \sum_{k=1}^{n}k^{6}&=& \frac{n(2n+1) (n+1)(3n^{4}+6n^{3}-3n+1)}{42} \label{sumas6}
  \\ \sum_{k=1}^{n}k^{7}&=& \frac{n^{2}(n+1)^{2}\left( 3n^{4}+6n^{3}-n^{2}-4n+2\right)}{24} \label{sumas7}
  % \\  \sum_{k=1}^{n}k^{8}&=&  \frac{n(2n+1)(n+1)(9n-n^{2}-15n^{3}+5n^{4}+15n^{5}+5n^{6}-3)}{90} \nonumber
\end{eqnarray}

\section[Sólidos de revolución en coordenadas paramétricas]{Sólidos de revolución en \texorpdfstring{\\ }{ } coordenadas paramétricas}

Los resultados de esta sección han sido estudiados muy posiblemente como aplicaciones de la integración en una variable,
algunos elementos como el concepto de \emph{momento} pueden consultarse particularmente en \cite{thomas_calculo_2015}. \index{Momento}
Los hemos utilizado para comprobar varios de los cálculos presentados en estas notas.

Consideremos una región $R$ acotada por las figuras de $f(x)$ y $g(x)$ con $x \in[a,b]$
con densidad constante $\rho$ que gira alrededor un una recta.

\begin{figure}[H]
  \centering
  \caption{La región gira sobre una recta exterior}
  \includegraphics[width=0.4\textwidth]{Fig/193.pdf} \hspace{5mm}
  \includegraphics[width=0.4\textwidth]{Fig/194.pdf}
  \fuentepropia
\end{figure}

\begin{teorema}[\bfseries teorema de Pappus para volúmenes] \index{Teorema de Pappus!volumen}
  Si una región plana girar alrededor de una recta en el plano, tal que la recta no corta el interior de la región, entonces el volumen $V$ del sólido generado es el área $A$ de la región por la distancia
  $2\pi r$ recorrida por su centroide durante la rotación. Si $r$ es la distancia entre el eje de rotación y
  el centroide, es decir $V=2\pi rA$.
\end{teorema}
\begin{proof}
  Verificaremos el resultado dos casos particulares.
  \begin{parrafoitemize}
  \item Denotaremos por $(\overline{x},\overline{y})$ a las coordenadas del centroide de $R$,
    su área está dada por \[A=\int_{a}^{b}\left(f(x)-g(x) \right) dx.\]
    Si $R$ gira alrededor de la recta $x=c$, entonces el volumen $V$ del sólido generado está dado por
    \[V=2\pi \int_{a}^{b}(x-c) \left( f(x) -g(x) \right) dx.\]
    Si $\mu_y$ denota el momento con respecto al eje $y$, es conocido que $\mu_y= \overline{x}\rho A$.
    El momento con respecto a la recta $L$, con ecuación $x=c$ está dado por
    \[\mu _{L}=\int_{a}^{b}\rho (x-c) \left( f(x)-g(x) \right) dx=\rho \int_{a}^{b}(x-c) \left(f(x) -g(x) \right) dx.\]
    Por otra parte,
    \[\mu_{L}=\rho\int_{a}^{b}x\left(f(x)-g(x)\right)dx-c\rho\int_{a}^{b}\left(f(x)-g(x)\right)dx=\mu_y-c\rho A.\]
    De las ecuaciones anteriores se sigue que $\mu_{L}=\rho \overline{x} A-c\rho A= \rho A (\overline{x}-c)$.
    Si $r$ es la distancia que recorre el centroide, $r=\overline{x}-c$ entonces, $V= 2 \pi r A$.
  \item Con las mismas consideraciones anteriores, si la región $R$ gira alrededor de la recta $y=c$,
    el volumen del sólido generado es
    \[V=\pi \int_{a}^{b}\left[ \left( f(x) -c\right) ^{2}-\left(g(x) -c\right) ^{2}\right] dx.\]
    Si $\mu_x$ denota el momento con respecto al eje $x$ es conocido que
    \[\mu_x= \frac{1}{2}\rho \int_{a}^{b}\left( [f(x)]^2 -[g(x)]^2\right)dx;\]
    por tanto $\mu_x= \rho\overline{y} A$. El momento con respecto a la recta $L$ con ecuación $y=c$, está dado por
    \begin{eqnarray*}
      \mu _{L}&=&\rho \int_{a}^{b}\left( f(x) -g(x) \right)\left( \frac{f(x) +g(x) }{2}-c\right) dx \\
      & =& \frac{1}{2}\rho \int_{a}^{b}\left( f(x) -g(x)\right) \left( f(x) +g(x) -2c\right) dx \\
      &=& \frac{1}{2}\rho \int_{a}^{b}\left[ \left( f(x) -c\right) ^{2}-\left(g(x) -c\right) ^{2}\right] dx
    \end{eqnarray*}
    Entonces, $\rho V = 2\pi \mu_L$; además,
    % $\mu_{L}=\frac{1}{2}\rho\int_{a}^{b} \left([f(x)]^2-[g(x)]^2\right)dx- c\rho\int_{a}^{b}\left(f(x) -g(x)\right)dx= \mu_x-c\rho A$.
    \begin{eqnarray*}\mu_{L} &=& \frac{1}{2}\rho\int_{a}^{b}\left([f(x)]^2-[g(x)]^2\right)dx-c\rho\int_{a}^{b}\left(f(x) -g(x)\right)dx. \\&=& \mu_x-c\rho A.
    \end{eqnarray*}
    Con las igualdades anteriores se sigue que
    \[\mu_{L}=\mu_x-c\rho A=\rho A \overline{y}-c\rho A= \rho A (\overline{y}-c).\]
    Si $r=\overline{y}-c$ es la distancia que recorre el centroide del eje de giro, entonces
    \[V=2\pi A (\overline{y}-c)=2\pi r A.\]
  \end{parrafoitemize}
\end{proof}

\subsubsection{Longitud de una curva en coordenadas paramétricas}

Para la curva $C$ parametrizada en la forma ${\bf r}(t)=\langle f(t),g(t)\rangle$ con $t\in [a,b]$,
consideremos dos puntos sobre la curva
$P_{i}=\left( x_{i},y_{i}\right) =\left( f\left( t_{i}\right) ,g\left(t_{i}\right) \right) $
y $P_{i+1}=\left( x_{i+1},y_{i+1}\right) =\left( f\left( t_{i}+\Delta t\right),g\left( t_{i}+\Delta t\right) \right) $
es posible establecer una aproximación a la longitud de arco.
\begin{eqnarray*}
  \overline{P_{i}P_{i+1}}&=& \sqrt{\left( f\left( t_{i}+\Delta t\right) -f\left( t_{i}\right) \right) ^{2}+\left( g\left( t_{i}+\Delta t\right) -g\left( t_{i}\right) \right) ^{2}} \\
  &=& \sqrt{\left( \frac{f\left( t_{i}+\Delta t\right) -f\left(t_{i}\right) }{\Delta t}\right) ^{2}+\left( \frac{g\left( t_{i}+\Delta t\right) -g\left( t_{i}\right) }{\Delta t}\right) ^{2}}\Delta t.
\end{eqnarray*}
La aproximación mejora si $\Delta t\to 0$ que equivale a que el número de rectángulos satisface $n \to \infty$.
Por lo tanto la longitud de la curva se calcula como
\begin{eqnarray*}
  L&=&\lim_{n\rightarrow \infty } \sum_{i=1}^{n}\sqrt{\left( \frac{%
    f\left( t_{i}+\Delta t\right) -f\left( t_{i}\right) }{\Delta t}\right)
    ^{2}+ \left( \frac{g\left( t_{i}+\Delta t\right) -g\left( t_{i}\right) }{%
  \Delta t}\right) ^{2}}\Delta t\\
  &=& \int_{a}^{b}\sqrt{\left(f^{\prime}(t)\right)^{2}+\left(g^{\prime}(t)\right)^{2}}dt=\int_{a}^{b}\Vert{\bf r}^\prime(t)\Vert dt.
\end{eqnarray*}
Esta expresión es la misma ecuación \ref{curvas}.
\subsubsection{Volumen de sólidos de revolución en coordenadas \\ paramétricas}

Si la curva $C$ parametrizada en la forma ${\bf r}(t)=\langle f(t),g(t)\rangle$ con $t\in [a,b]$, gira alrededor del eje $x$
entonces el volumen del sólido obtenido se puede determinar de la siguiente forma.

\begin{figure}[H]
  \centering
  \caption{Sólido de revolución alrededor del eje $x$}
  \includegraphics[width=0.4\textwidth]{Fig/195.pdf} \hspace{5mm}
  \includegraphics[width=0.4\textwidth]{Fig/196.pdf}
  \fuentepropia
\end{figure}

Para un rectángulo de altura $g\left( t_{i}^{\ast }\right) $ y
base $x_{i+1}-x_{i}$ se sigue que
\[x_{i+1}-x_{i}=f\left( t_{i}+\Delta t\right) -f\left( t_{i}\right) =\left(
\frac{f\left( t_{i}+\Delta t\right) -f\left( t_{i}\right) }{\Delta t}\right)\Delta t.\]
El volumen del solido generado por la rotación del $i-$ésimo rectángulo está dado por
\begin{eqnarray*}
  V_{i} &=&\pi \left( g\left( t_{i}^{\ast }\right) \right) ^{2}(f\left( t_{i}+\Delta t\right) -f\left( t_{i}\right));
  %\\ & =& \pi \left( g\left( t_{i}^{\ast }\right) \right) ^{2}\left( \frac{f\left( t_{i}+\Delta t\right) -f\left( t_{i}\right) }{\Delta t}\right)\Delta t
\end{eqnarray*}
entonces,

\begin{eqnarray*}
  V &=& \lim_{n\rightarrow \infty }\left( \sum_{i=1}^{n}\pi \left(g\left( t_{i}^{\ast }\right) \right) ^{2}\left( \frac{f\left( t_{i}+\Delta t\right) -f\left( t_{i}\right) }{\Delta t}\right) \Delta t\right)\\ &=& \int_{a}^{b}\pi \left[ g\left( t\right) \right] ^{2}f^{\prime }(t) dt.
\end{eqnarray*}

\begin{figure}[H]
  \centering
  \caption{Sólido de revolución}
  \includegraphics[width=0.4\textwidth]{Fig/197.pdf} \hspace{5mm}
  \includegraphics[width=0.4\textwidth]{Fig/198.pdf}
  \fuentepropia
\end{figure}

Esta expresión comúnmente se escribe de manera simplificada como
\[V=\int_{a}^{b}\pi \left[ y\left( t\right) \right] ^{2}x^{\prime }(t) dt\]
Al girar alrededor del eje $y$ la expresión correspondiente es
\[V=\int_{a}^{b}\pi \left[ x\left( t\right) \right] ^{2}y^{\prime }(t) dt.\]
Otro resultado importante es el siguiente:

\begin{teorema}[\bfseries Teorema de Pappus para áreas superficiales] \index{Teorema de Pappus!área superficial}\label{teoremadepappus}
  Si un arco de una curva plana suave se hace girar alrededor de una recta
  en el plano, tal que la recta no corta al interior del arco, entonces el
  área $A$ de la superficie generada por el arco es igual a la longitud del arco $L$ por la
  distancia recorrida por el centroide del arco $2\pi r$ durante la rotación, es decir $A=2\pi rL$.
\end{teorema}

%\clearpage

\begin{proof}
  Para una curva en el plano $xy$ descrita en forma paramé- \\ trica
$\left\langle f(t) ,g(t) \right\rangle$ con $t\in\lbrack a,b]$. La longitud $L$ de la curva está dada
\[L=\int_{a}^{b}\sqrt{\left( g^{\prime }(t)\right) ^{2}+\left( f^{\prime }(t) \right) ^{2}}dt.\]
Si la curva gira alrededor del eje $y$, entonces la función
\[{\bf r}\left(u,t\right) =\left\langle f(t) \cos u,g(t) ,f(t) \sin u\right\rangle,\]
con $u\in \lbrack 0,2\pi ]$
describe la superficie de revolución.
Con ${\bf r}$ hallamos los siguientes vectores
\begin{eqnarray*}
{\bf r}_{u} &=& \left\langle -f(t) \sin u,0,f(t) \cos u\right\rangle \\
{\bf r}_{t}&=&\left\langle f^{\prime }(t) \cos u,g^{\prime }(t) ,f^{\prime }(t) \sin u\right\rangle \\
{\bf r}_{u}\times {\bf r}_{t}&=& \left\langle-f(t)g^{\prime}(t)\cos u,f(t)f^{\prime}(t),-f(t)g^{\prime}(t)\sin u\right\rangle;
\end{eqnarray*}
entonces:
\begin{eqnarray*}
\left\Vert {\bf r}_{u}\times {\bf r}_{t}\right\Vert &=& \left\vert f(t) \right\vert \left\Vert\left\langle -g^{\prime }(t) \cos u,f^{\prime }(t),-g^{\prime }(t) \sin u\right\rangle \right\Vert\\
&=& \left\vert f(t) \right\vert \left( \left( -g^{\prime }(t)\right) ^{2}\cos ^{2}u+\left( f^{\prime }(t) \right) ^{2}+\left(-g^{\prime }(t) \right) ^{2}\sin ^{2}u\right) \\
&=& \left\vert f(t) \right\vert \sqrt{\left( g^{\prime }(t) \right) ^{2}+\left( f^{\prime }(t) \right) ^{2}}.
\end{eqnarray*}
Por lo tanto, el área de la superficie está dada por
%
\begin{eqnarray*}
A &=& \int_{a}^{b}\left( \int_{0}^{2\pi }\left\vert f(t)\right\vert \sqrt{\left( g^{\prime }(t) \right) ^{2}+\left(f^{\prime }(t) \right) ^{2}}du\right) dt \\
&=& 2\pi\int_{a}^{b}\left\vert f(t) \right\vert \sqrt{\left( g^{\prime}(t) \right) ^{2}+\left( f^{\prime }(t) \right) ^{2}}dt
\end{eqnarray*}

Si el giro se hace alrededor del eje $x$, la fórmula respectiva es
\[A=2\pi \int_{a}^{b}\left\vert g(t) \right\vert \sqrt{\left(g^{\prime }(t) \right) ^{2}+\left( f^{\prime }(t)\right) ^{2}}dt.\]
La masa $M$ de un alambre cuya forma la determina el arco en mención está dada por
\[M=\rho \int_{a}^{b}\sqrt{\left( g^{\prime }(t)\right) ^{2}+\left( f^{\prime }(t) \right) ^{2}}dt.\]
Supongamos que $f(t)\geq 0$ y  $g(t)\geq 0$, entonces las coordenadas del centroide de la curva están dadas por
\begin{eqnarray}
\overline{x} %&=&\frac{1}{M}\int_{a}^{b}\rho f(t) \sqrt{\left( g^{\prime }(t) \right) ^{2}+\left( f^{\prime }(t) \right) ^{2}}dt\nonumber \\
&=&\frac{\rho }{M}\int_{a}^{b}f(t) \sqrt{\left( g^{\prime }(t) \right) ^{2}+\left( f^{\prime }(t) \right) ^{2}}dt \label{pappus1}
\end{eqnarray}
\begin{eqnarray}
\overline{y} %&=& \frac{1}{M}\int_{a}^{b}\rho g(t) \sqrt{\left(g^{\prime }(t) \right) ^{2}+\left( f^{\prime }(t)\right)^{2}}dt \nonumber  \\
&=& \frac{\rho }{M}\int_{a}^{b}g(t) \sqrt{\left(g^{\prime }(t) \right) ^{2}+\left( f^{\prime }(t)\right) ^{2}}dt\label{pappus2}
\end{eqnarray}
Con lo anterior es posible establecer la igualdad
\begin{eqnarray*}
A &=& 2\pi \int_{a}^{b}f(t) \sqrt{\left( g^{\prime }(t) \right) ^{2}+\left( f^{\prime }(t) \right) ^{2}}dt \\
&=& 2\pi\left( \frac{M}{\rho }\right) \overline{x}=2\pi \left( \frac{1}{\rho }\right) \overline{x}\rho \int_{a}^{b}\sqrt{\left( g^{\prime }(t)\right) ^{2}+\left( f^{\prime }(t) \right) ^{2}}dt \\
&=& 2\pi \overline{x}\int_{a}^{b}\sqrt{\left( g^{\prime }(t)\right) ^{2}+\left( f^{\prime }(t) \right) ^{2}}dtA=2\pi \overline{x}L.
\end{eqnarray*}
Si el giro se da alrededor del eje $y$, de la fórmula \ref{pappus2} se sigue que
\begin{eqnarray*}
A &=& 2\pi \int_{a}^{b}g(t)\sqrt{\left( g^{\prime }(t) \right) ^{2}+\left( f^{\prime }(t) \right) ^{2}}dt \\
&=& 2\pi\left( \frac{M}{\rho }\right) \overline{y}=2\pi \left( \frac{1}{\rho }\right) \overline{y}\rho \int_{a}^{b}\sqrt{\left( g^{\prime }(t)\right) ^{2}+\left( f^{\prime }(t) \right) ^{2}}dt \\
&=& 2\pi \overline{y}\int_{a}^{b}\sqrt{\left(g^{\prime}(t)\right)^{2}+\left(f^{\prime}(t)\right)^{2}}dtA=2\pi \overline{y}L.
\end{eqnarray*}
\end{proof}

\begin{figure}[H]
\centering
\caption{Superficie de revolución}
\includegraphics[width=0.8\textwidth]{Fig/199.pdf}
\fuentepropia
\end{figure}

\begin{ejer}
La curva ${\bf r}(t)=\big\langle 3+(2-\sin (t))\cos (t),0,3\sin (t)(2+\cos
(2t))\big \rangle $, con $t\in \lbrack 0,2\pi ]$ gira alrededor del eje $y$, generando una superficie de revolución.

Hallar el área de la superficie genrada y el volumen del slido acotado por esta región.

En el siguiente enlace
\href{https://www.math3d.org/zbgohcryUD}{https://www.math3d.org/zbgohcryUD}, puede encontrase esta construcción.

\begin{figure}[H]
\centering
\caption{Superficie de revolución}
\includegraphics[width=0.4\textwidth]{Img/T_74a.pdf} \hspace{5mm}
\includegraphics[width=0.4\textwidth]{Img/T_74b.pdf}

\includegraphics[width=0.4\textwidth]{Img/T_74c.pdf}
\fuentepropia
\end{figure}

La superficie de revolución obtenida se parametriza por
{\small
\begin{equation*}
\begin{split}
  {\bf r}(u,v)&=\big \langle (3+(2-\sin (u))\cos (u))\cos (v),(3+(2-\sin (u))\cos(u))\sin (v), \\
  &\qquad 3\sin (u)(2+\cos (2u))\big \rangle
\end{split}
\end{equation*}
}

con $u\in \lbrack 0,2\pi ]$, $v\in \lbrack 0,2\pi]; $ pues,
el área y el volumen requeridos pueden hallarse con los métodos expuestos antes.
\end{ejer}

\section[Coordenadas curvilíneas ortogonales]{Coordenadas curvilíneas \texorpdfstring{\\ }{ } ortogonales}

Un cambio de coordenadas cartesianas $(x,y,z)$ a otras diferentes $(u,v,w)$ es una transformación
\emph{uno a uno} que consta de tres funciones con derivadas parciales continuas
$x = x(u,v,w)$, $y=y(u,v,w)$, $z=z(u,v,w)$,
cuya transformación inversa $u=u(x,y,z)$, $v=v(x,y,z)$, $w = w(x,y,z)$, existe.
Esto es posible siempre que \index{Coordenadas curvilíneas!ortogonales}

\[\frac{\partial (x,y,z)}{\partial (u,v,w)}= \left|
\begin{array}{ccc}
\frac{\partial x}{\partial u} & \frac{\partial x}{\partial v} & \frac{\partial x}{\partial w}\\[0.1cm]
\frac{\partial y}{\partial u} & \frac{\partial y}{\partial v} & \frac{\partial y}{\partial w}\\[0.1cm]
\frac{\partial z}{\partial u} & \frac{\partial z}{\partial v} & \frac{\partial z}{\partial w}
\end{array} \right| \neq 0 \]

Si ${\bf r}(u,v,w)=\la x(u,v,w),y(u,v,w),z(u,v,w)  \ra$
denota el vector posición,
una \emph{superficie coordenada} consta de todos los puntos en los cuales las funciones  $u$, $v$ o $w$ toman valores constantes. La intersección de dos superficies coordenadas se denomina \emph{curva coordenada}.

El sistema se denomina \emph{ortogonal} si las tangentes a las curvas coordenadas en todo punto de intersección,
son ortogonales. Es decir, las coordenadas $u,v,w$ forman un sistema de \emph{coordenadas curvilíneas ortogonales} si el conjunto

\[\left \{\frac{\partial{\bf r}}{\partial u},\frac{\partial{\bf r}}{\partial v},\frac{\partial{\bf r}}{\partial w}\right \}\]

es una base ortogonal y orientada positivamente. Los factores de escala se determinan mediante las siguientes expresiones
\[h_u= \left \Vert \frac{\partial{\bf r}}{\partial u}\right \Vert,\, h_v=\left\Vert\frac{\partial{\bf r}}{\partial v}\right\Vert,\,  h_w=\left\Vert\frac{\partial{\bf r}}{\partial w}\right\Vert\]
Al definir
\[\ei{u}= \frac{1}{h_u}\frac{\partial{\bf r}}{\partial u},\, \ei{v}=\frac{1}{h_v}\frac{\partial{\bf r}}{\partial v},\, \ei{w}=\frac{1}{h_w}\frac{\partial{\bf r}}{\partial w}\]

se sigue que el conjunto  $\{ \ei{u},\ei{v},\ei{w} \}$
es una base ortonormal y orientada positivamente.
El gradiente de un campo escalar $f$ en $\mathbb{R}^{3}$ está dado por
\[\nabla f=\la \frac{\partial f}{\partial x},\frac{\partial f}{\partial y},
\frac{\partial f}{\partial z}\ra,\]
en coordenadas curvilíneas ortogonales la expresión correspondiente es
\[\nabla f= \frac{1}{h_{u}}\frac{\partial f}{\partial u}\ei{u}+\frac{1}{h_{v}}\frac{\partial f}{\partial v}\ei{v}+\frac{1}{h_{w}}\frac{\partial f}{\partial w}\ei{w}.\]
La divergencia de un campo vectorial ${\bf F}\left(x,y,z\right)$
% $=\big\langle P(x,y,z),Q(x,y,z),R(x,y,z)\big\rangle$
% es el campo escalar ${\rm div}{\bf F}=\frac{\partial P}{\partial x}+\frac{\partial Q}{\partial y}+ \frac{\partial R}{\partial z}$
en coordenadas curvilíneas ortogonales la divergencia está dada por
\index{Coordenadas curvilíneas!divergencia}
\begin{multline*}
{\rm div}{\bf F}=\frac{1}{h_{u}h_{v}h_{w}}\bigg(
\frac{\partial}{\partial u} \left(h_{v}h_{w}{\bf F}(u,v,w)\right)
+\frac{\partial}{\partial v}\left(h_{u}h_{w}{\bf F}(u,v,w) \right) \\
+\frac{\partial}{\partial w}\left(h_{u}h_{v}{\bf F}(u,v,w) \right) \bigg).
\end{multline*}
%El rotacional del campo vectorial ${\bf F}$ en coordenadas cartesianas, se puede calcular por la expresión \index{Coordenadas curvilíneas!rotacional}
%${\rm rot} {\bf F}=\left\vert \begin{array}{ccc} \vi & \vj & \vk \\[0.1cm] \frac{\partial }{\partial x} & \frac{\partial }{\partial y} & \frac{\partial  }{\partial z} \\[0.1cm]   P(x,y,z)  & Q(x,y,z)  & R(x,y,z) \end{array}\right\vert $
El rotacional en coordenadas curvilíneas  de ${\bf F}$ está determinado por
{\small
\[{\rm rot}{\bf F}=\frac{1}{h_{u}h_{v}h_{w}}\left\vert
\begin{array}{ccc}
h_{u}\ei{u} & h_{v}\ei{v} & h_{w}\ei{w} \\[0.15cm]
\frac{\partial }{\partial u} & \frac{\partial }{\partial v} & \frac{\partial
}{\partial w} \\[0.15cm]
h_{u}P\left( u,v,w\right)  & h_{v}Q\left( u,v,w\right)  & h_{w}R\left(
u,v,w\right)
\end{array} \right\vert\]
}

Una exposición más profunda de estos temas puede consultarse en \\ \cite{poljanin_handbook_2007}

Presentamos algunos tipos de sistemas de coordenadas curvilíneas.
\subsection{Coordenadas curvilíneas en el plano} \index{Coordenadas curvilíneas}
\subsubsection{Coordenadas polares}
Este sistema de coordenadas es posiblemente muy conocido por el lector, la relación con las coordenadas
cartesianas es
\[x = r \cos \theta, \hskip0.15cm y=r \sin \theta,\hskip0.15cm r\in [ 0,\infty),\,\theta \in [0 ,2\pi ).\]

De las ecuaciones $x^2+y^2=r^2$
y $\tan \theta = y/x$
es posibles establecer que para valores constantes de $r$
las curvas coordenadas son circunferencias, y si $\theta$ es constante,
se generan semirectas. El punto central ($r=0$) se denomina \emph{polo,}
y la semirrecta $\theta=0$ se denomina \emph{eje polar.}

\begin{figure}[H]
\centering
\caption{Plano polar}
\includegraphics[width=0.5\textwidth]{Fig/200.pdf}
\fuentepropia
\end{figure}

\subsubsection{Coordenadas bipolares} \index{Coordenadas!bipolares}
Para este sistema de coordenadas en el plano se define
\begin{equation}
x =\frac{a\sinh v}{\cosh v-\cos u} ,\hskip0.5cm y=\frac{a\sin u}{\cosh v-\cos u}.
\end{equation}
Con $u\in \lbrack 0,2\pi )$, $v\in (-\infty ,\infty )$.

Presentamos algunas curvas en coordenadas bipolares, en este caso las curvas coordenadas poseen dos polos ubicados en $(-a,0)$ y $(a,0)$.
Esta afirmación es consecuencia de los siguientes resultados
\[\lim_{v\rightarrow \infty }x=a,\, \lim_{v\rightarrow -\infty }x= -a\]
\[\lim_{v\rightarrow\infty}y=\lim_{v\rightarrow -\infty }y=0.\]

\begin{figure}[H]
\centering
\caption{Coordenadas bipolares}
\includegraphics[width=0.5\textwidth]{Fig/201.pdf}
\fuentepropia
\end{figure}

\begin{prop}
En coordenadas bipolares con valores constantes de $u$ y $v$ representan c\'{\i}rculos en el plano $xy$.
\end{prop}

\begin{proof}
Para tal fin mostraremos las siguientes igualdades:
$x^2+\left( y-a\cot u\right) ^2=a^2\csc ^2u;\hskip0.2cm  \left( x-a\coth v\right) ^2+y^2=a^2{\rm csch}^2v$.
Para la primera de estas expresiones
{\small
\begin{multline}
\label{bipolares1}
x^2+\left( y-a\cot u\right) ^2 =  \left( \frac{a\sinh v}{\cosh v-\cos u}\right) ^2+\left( \frac{a\sin u}{\cosh v-\cos u}-\frac{a\cos u}{\sin u}\right) ^2  \\
=  \frac{a^2\left( \sinh ^2v\sin ^2u+\left( \cos u\cosh v-1\right) ^2\right)}{\left( \cosh v-\cos u\right) ^2\sin ^2u}   \\
=  \frac{a^2\left( \cosh v-\cos u\right) ^2 }{\left( \cosh v-\cos u\right) ^2\sin ^2u} = %\frac{a^2}{\sin ^2u} \\ =
a^2{\rm csch}^2u
\end{multline}
}
% = \frac{a^2\left(\left( \sinh v\sin u\right) ^2+\left( \sin ^2u-\left( \cosh v-\cos u\right) \cos u\right) ^2\right) }{\left( \cosh v-\cos u\right) ^2\sin ^2u}
La otra igualdad se trabaja de manera similar
\begin{multline} \label{bipolares2}
(x-a\coth v) ^2+y^2 = \left( \frac{a\sinh v}{\cosh v-\cos u}-\frac{a\cosh v}{\sinh v}\right) ^2+ \\ \left( \frac{a\sin u}{\cosh v-\cos u}\right) ^2
= \frac{a^2\left( \cos u-\cosh v\right) ^2}{\left( \cosh v-\cos u\right) ^2\sinh ^2v}  %\\ = \frac{a^2}{\sinh ^2v}
= a^2\csch^2v
\end{multline}
\end{proof}
% = a^2\left( \left( \frac{\sinh ^2v-\left( \cosh v-\cos u\right) \cosh v}{ \left( \cosh v-\cos u\right) \sinh v}\right) ^2+\left( \frac{\sin u}{\cosh v-\cos u}\right) ^2\right) = \frac{a^2\left( \left( \sinh ^2v-\left( \cosh v-\cos u\right) \cosh v\right) ^2+\sin ^2u\sinh^2v\right) }{\left( \cosh v-\cos u\right) ^2\sinh ^2v} = \frac{a^2\left( \left( \cos u\cosh v-1\right) ^2+\sin ^2u\sinh ^2v\right)}{\left( \cosh v-\cos u\right) ^2\sinh ^2v}

\subsubsection{Coordenadas elípticas} \index{Coordenadas!elípticas}
La relación entre coordenadas cartesianas y coordenadas elípticas es
\begin{equation}
x = a \cosh u\cos v,\hskip0.5cm y = a \sinh u \sin v.
\end{equation}
En donde $u \in \mathbb{R}$ y $v \in [0, 2\pi]$.

Las líneas coordenadas son elípses e hipérbolas confocales (con los mismos focos). La distancia focal es $a$. \index{Confocales}

\begin{figure}[H]
\centering
\caption{Curvas confocales}
\includegraphics[width=0.5\textwidth]{Fig/202.pdf}
\fuentepropia
\end{figure}

Se puede mostrar que la obtención de dichas trazas es consecuencia de
las siguientes igualdades.
\begin{equation}
\frac{x^2}{a^2\cosh ^2u}+\frac{y^2}{a^2\sinh ^2u} =  1 \label{elipticas1}
\end{equation}
\begin{equation}
\frac{x^2}{a^2\cos ^2v}-\frac{y^2}{a^2\sin ^2v} = 1  \label{elipticas2}
\end{equation}

\subsubsection{Coordenadas curvilíneas en el espacio} \index{Coordenadas!curvilíneas en el espacio}
Algunos ejemplos de los presentados en esta sección son una extensión de
las respectivas coordenadas en el plano.

\paragraph*{Coordenadas cilíndricas elípticas}
\index{Coordenadas!cilíndricas elípticas}
En este sistema de coodenadas $u$ y $v$ representan cilindros hiperbólicos confocales y cilindros elípticos confocales.
La conversión de coordenadas cilíndricas elípticas a rectangulares es la siguiente.
\begin{equation}
x = a \cosh u\cos v,\hskip0.5cm y = a \sinh u \sin v, \hskip0.5cm z=z
\end{equation}
Con $u\in \lbrack 0,\infty )$,  $v\in \lbrack 0,2\pi )$ y  $z\in (-\infty ,\infty )$.
A partir de la ecuación \ref{elipticas1} se evidencia que para valores constantes
de $u$ la superficie determinada es un cilindro elíptico;
mientras que la expresión \ref{elipticas2} representa en el espacio un cilíndro hiperbólido.

\begin{figure}[H]
\centering
\caption{Superficies constantes en coordenadas cilíndricas elípticas}
\includegraphics[width=0.4\textwidth]{Img/T_63a.pdf} \hspace{5mm}
\includegraphics[width=0.4\textwidth]{Img/T_63b.pdf}
\fuentepropia
\end{figure}

Estas figuras se consiguen con las siguientes instrucciones.

\begin{tcolorbox}[title=\bf Instrucción \verb"Mathematica"]\begin{Verbatim}[formatcom=\letraverbatim]
{x,y,u}={Cosh[u]*Cos[v],Sinh[u]*Sin[v],2};
A1=ParametricPlot3D[{x,y,z},{v,0,2*Pi},{z,-3,3}];
v=5Pi/4;
A2=ParametricPlot3D[{x,y,z},{u,-3,3},{z,-3,3}];
z=1;
A3=ParametricPlot3D[{x,y,z},{u,0,2*Pi},{v,-3,3}];
Show[A2,A3,PlotRange->{{-4,4},{-4,4},{-3,3}}]
\end{Verbatim}
\end{tcolorbox}

\subsubsection{Coordenadas cil\'{\i}ndricas parabólicas}  \index{Coordenadas!cilíndricas parabólicas}
La relación con las coordenadas cartesianas y el sistema de coordenadas cil\'{\i}ndricas parabólicas son las siguientes:
\begin{equation}
x = \frac{1}{2}\left( u^2-v^2\right),\, y = uv, \, z=z \
\end{equation}
Con $u\in \lbrack 0,\infty )$,  $v\in \lbrack 0,\infty )$ y  $z\in (-\infty,\infty)$.

Las trazas en el plano $xy$ son parábolas, pues para valores constantes de $u$ la ecuación $y^2=u^2\left( u^2-2x\right)$,
representa una parábola que abre hacia la negativa del eje $x$; y para cuando $v$ es constante la ecuación
$y^2=v^2\left( v^2+2x\right)$, caracteriza una parábola abre hacia la parte positiva del eje $x$.

\begin{figure}[H]
\centering
\caption{Coordenadas bipolares}
\includegraphics[width=0.5\textwidth]{Fig/203.pdf}
\fuentepropia
\end{figure}

Es por esto que en el espacio  las superficies constantes son cilindros parabólicos.
Las consideraciones anteriores también se pueden evidenciar con las siguientes igualdades:
{\small
\[u^2\left( u^2-2x\right) =u^2\left( u^2-2\left( \frac{1}{2}\left(
u^2-v^2\right) \right) \right) = u^2v^2=y^2\]
\[v^2\left( v^2+2x\right) =v^2\left( v^2+2\left( \frac{1}{2}\left(
u^2-v^2\right) \right) \right) = u^2v^2=y^2\]
}

\begin{figure}[H]
\captionof{figure}{\footnotesize Superficies constantes en coordenadas cilíndricas parabólicas.}
\centering
\includegraphics[width=4.5cm,height=3.5cm]{Img/T_66a.pdf}
\includegraphics[width=4cm,height=3.5cm]{Img/T_66b.pdf}
\fuentepropia
\end{figure}

\begin{tcolorbox}[title=\bf Instrucción \verb"Mathematica"]\begin{Verbatim}[formatcom=\letraverbatim]
{x,y}={(u^2-v^2)/2,u*v};
u=2;
A1=ParametricPlot3D[{x,y,z},{v,-3,3},{z,-2,2}];
v=2;
A2=ParametricPlot3D[{x,y,z},{u,-3,3},{z,-2,2}];
z=1;
A3=ParametricPlot3D[{x,y,z},{u,-3,3},{v,-3,3}];
Show[A1,A2,A3,PlotRange->{{-4,4},{-5,5},{-2,2}}]
\end{Verbatim}
\end{tcolorbox}

\subsubsection{Coordenadas parabólicas} \index{Coordenadas!parabólicas}
Según \cite{hernandez_nolasco_familias_2012}, este sistema de coordenadas se genera por la rotación alrededor del $z$ de dos parábolas bidimensionales;
la relación con las coordenadas cartesianas es la siguiente
\begin{equation}
x = uv\cos \theta,\hskip0.5cm y = uv\sin \theta, \hskip0.5cm z = \frac{1}{2}\left( u^2-v^2\right)
\end{equation}

Con $u\in [ 0,\infty )$, $v\in [ 0,\infty )$, $\theta \in [0,2\pi )$.

De las relaciones $x^2+y^2=u^2v^2=u^2(u^2-2z)$ y $x^2+y^2=u^2v^2=v^2(2z-v^2)$ se observa que para
valores constantes $u$ la superficie obtenida es un paraboloide que abre hacia la parte negativa del eje $z$, y para valores constantes de $v$ también se obtiene paraboloides, solo que estos abren hacia a parte positiva del eje $z$. Además, $\tan \theta=y/x$ permite concluir que valores constante de $\theta$ producen semiplanos.

\begin{figure}[H]
\centering
\caption{Superficies constantes en coordenadas parabólicas}
\includegraphics[width=0.4\textwidth]{Img/T_64a.pdf} \hspace{5mm}
\includegraphics[width=0.4\textwidth]{Img/T_64b.pdf}
\fuentepropia
\end{figure}

\begin{tcolorbox}[title=\bf Instrucción \verb"Mathematica"]\begin{Verbatim}[formatcom=\letraverbatim]
Clear[u,v,\[Theta]];
{x,y,z}={u*v*Cos[\[Theta]],u*v*Sin[\[Theta]],
(u^2-v^2)/2};
u=2;
A1=ParametricPlot3D[{x,y,z},{v,0,2},{\[Theta],0,2*Pi}]
v=2;
A2=ParametricPlot3D[{x,y,z},{u,0,2},{\[Theta],0,2*Pi}]
\[Theta]=Pi/3;
A3=ParametricPlot3D[{x,y,z},{u,0,2},{v,0,2},
PlotStyle->Yellow];
Show[A1,A3,PlotRange->{{-4,4},{-4,4},{-3,3}}]
\end{Verbatim}
\end{tcolorbox}

Otras relaciones que pueden resultar útiles son:
\[x^2+y^2+z^2=\left( uv\cos \theta \right) ^2+\left( uv\sin \theta
\right) ^2+\left( \frac{1}{2}\left( u^2-v^2\right) \right) ^2= \frac{1}{4}\left( u^2+v^2\right) ^2,\]
de la cual se sigue que
\begin{eqnarray*}
\sqrt{x^2+y^2+z^2}&=&\frac{1}{2}\left( u^2+v^2\right) \\
\sqrt{x^2+y^2+z^2}+z &=& \frac{1}{2}\left( u^2+v^2\right) +\frac{1}{2}\left( u^2-v^2\right) = u^2\\
\sqrt{x^2+y^2+z^2}-z &=& \frac{1}{2}\left( u^2+v^2\right) -\frac{1}{2}\left( u^2-v^2\right) = v^2
\end{eqnarray*}

\subsubsection{Coordenadas cilíndricas bipolares} \index{Coordenadas!cilíndricas bipolares}
Este sistema de coordenadas se relaciona con las coordenadas cartesianas de la siguiente manera.

\begin{equation}
x = \frac{a\sinh v}{\cosh v-\cos u},\hskip0.5cm y =\frac{a\sin u}{\cosh v-\cos u}, \hskip0.5cm z = z
\end{equation}
En donde $u\in \lbrack 0,2\pi )$, $v\in (-\infty ,\infty )$, $z\in (-\infty,\infty )$.
Utilizando las expresiones \ref{bipolares1} y \ref{bipolares2} se concluyen que en
coordenadas cilíndricas bipolares valores constantes de $u$ y $v$ generan cilindros circulares con $z$ como eje de simetría.

Algunas de estas figuras se consiguen con las siguientes instrucciones:

\begin{tcolorbox}[title=\bf Instrucción \verb"Mathematica"]\begin{Verbatim}[formatcom=\letraverbatim]
x=Sinh[v]/(Cosh[v]-Cos[u]);
y=Sin[u]/(Cosh[v]-Cos[u]);
u=Pi/6;
A1=ParametricPlot3D[{x,y,z},{v,-4,4},{z,-2,2},
MeshStyle->Black,PlotStyle->Orange,Mesh->10,
PlotPoints->50];
v=1/2;
A2=ParametricPlot3D[{x,y,z},{u,0,2*Pi},{z,-2,2},
MeshStyle->Gray,PlotPoints->50,PlotStyle->Cyan];
z=1;
A3=ParametricPlot3D[{x,y,z},{u,0,2*Pi},{v,-2,2},
MeshStyle->Gray,PlotPoints->50,PlotStyle->Orange];
Show[A1,A2,A3,PlotRange->{{-3,4},{-3,4},{-2,2}}]
\end{Verbatim}
\end{tcolorbox}

Cuyo producto es el siguiente \ref{fig:cilindricas_bipolares}

\begin{figure}[H]
\centering
\caption{Superficies constantes en coordenadas cilíndricas bipolares}
\label{fig:cilindricas_bipolares}
\includegraphics[width=0.4\textwidth]{Img/T_65a.pdf} \hspace{5mm}
\includegraphics[width=0.4\textwidth]{Img/T_65b.pdf}
\fuentepropia
\end{figure}

\begin{ejer}
Mostrar las siguientes igualdades
{\small
\[\tan u=\frac{2ay}{x^{2}+y^{2}-a^{2}}, \hskip0.15cm
\tanh v=\frac{2ax}{x^{2}+y^{2}+a^{2}},\hskip0.15cm
v=\frac{1}{2}\ln \left( \frac{\left(x+a\right)^{2}+y^{2}}{\left( x-a\right)^{2}+y^{2}}\right)\!.\]
}
\end{ejer}

\subsubsection{Coordenadas toroidales} \index{Coordenadas!toroidales}
\paragraph*{Construcción de un toro}

En el plano $yz$, la circunferencia de radio $r$ y centro en el punto $(a,0)$
con $a>r$ posee ecuación $\left( y-a\right) ^{2}+z^{2}=r^{2}$.
Esta curva se puede parametrizar en la forma $y=a+r\cos \theta $, $z=r\sin \theta $,
con $\theta \in [0,2\pi]$.

\begin{figure}[H]
\centering
\includegraphics[width=0.3\linewidth]{Fig/204.pdf}
\includegraphics[width=0.3\linewidth]{Fig/205.pdf}
%\caption{[falta título]}
\fuentepropia
\end{figure}

Si un punto sobre esta curva gira alrededor del eje $z$ describe una circunferencia, en
cuyo radio es el valor que toma $y$, la coordenada $z$ es la misma para
todos los puntos de dicha circunferencia.
La superficie generada por el giro de la circunferencia sobre el eje $z$ se denomina  un \emph{toro}
y se determina por medio de las expresiones \index{Toro}
\[x = \left( a+r\cos \theta \right) \cos u,\,\, y = \left( a+r\cos \theta \right) \sin u,\,\, z = r\sin \theta,\]
con $\theta \in \lbrack 0,2\pi )$, $u\in \lbrack 0,2\pi)$.
De estas ecuaciones se obtienen las siguientes relaciones:
\begin{eqnarray*}
x^{2}+y^{2}+z^{2} &=& \left( \left(a+r\cos\theta \right)\cos u\right)^{2}+
\left(\left(a+r\cos\theta\right)\sin u\right) ^{2}+r^{2}\sin^{2}\theta \\ &=& a^{2}+2ar\cos \theta +r^{2}.
\end{eqnarray*}
Además,
$\sqrt{x^{2}+y^{2}}=a+r\cos \theta \rightarrow r\cos \theta =\sqrt{x^{2}+y^{2}}-a$.
Por lo tanto,
\begin{equation} \label{toroidales3}
x^{2}+y^{2}+z^{2}+(a^{2}-r^{2})= 2a\sqrt{x^{2}+y^{2}}
\end{equation}
\paragraph*{Coordenadas toroidales} \index{Coordenadas!toroidales}
Este sistema de coordenadas se obtiene cuando se hace girar alrededor del eje $z$
las coordenadas bipolares en dos dimensiones.

% De acuerdo a\footnote{\verb"https://rmf.smf.mx/pdf/rmf/40/5/40_5_805.pdf"}
Las superficies correspondientes a valores constantes de $v$ corresponden a \emph{toros con su
corazón circular de radio $a$ en el plano ecuatorial,} cuando $u$ es constante se generan esferas ortogonales
a los toros cuyo centro está en el eje polar, y con $\phi$ constante se obtienen semiplanos.
Las siguientes ecuaciones determinan la relación entre coordenadas cartesianas y coordenadas toroidales.
\begin{equation}
x = \frac{a\sinh v\cos \phi }{\cosh v-\cos u},\hskip0.5cm y =\frac{a\sinh v\sin \phi }{\cosh v-\cos u}, \hskip0.5cm z = \frac{a\sin u}{\cosh v-\cos u}
\end{equation}

Con $u\in [0,2\pi )$, $v\in [0,\infty )$ y $\phi \in [0,2\pi )$.
Si $\phi $ es constante, es claro que la expresión $\tan \phi =y/x$ representa un semiplano.

Ahora se mostraran las otras afirmaciones hechas al inicio.

\begin{prop}
En coordenadas toroidales se determina un toro con valores constantes de $v$.
\end{prop}

\begin{proof}
De la definición se obtienen algunas identidades:
{\small
\begin{multline}
x^2+y^2 = \left( \frac{a\sinh v\cos \phi }{\cosh v-\cos u}\right)^2+\left( \frac{a\sinh v\sin \phi }{\cosh v-\cos u}\right) ^2\\
=  \frac{a^2 \sinh ^2v}{\left(\cosh v-\cos u\right) ^2}. \label{toroidales1}
\end{multline}
}

Lo cual implica que {\small $\sqrt{x^2+y^2}=\dfrac{a\sinh v}{\cosh v-\cos u}$.
}

Otra equivalencia es
\begin{eqnarray}
x^2+y^2+z^2+a^2 &=&  \frac{a^2 \sinh ^2v}{\left(\cosh v-\cos u\right) ^2} +\left( \frac{a\sin u}{\cosh v-\cos u}\right) ^2+a^2 \nonumber \\
% &=& \frac{a^2\left(\sinh^2 v+\sin^2 u+(\cosh v-\cos u)^2\right)}{\left(\cosh v-\cos u\right)^2}\nonumber\\
% &=& \frac{a^2\left(\sinh ^{2}v+\cosh ^{2}v-2\cosh v\cos u+1 \right)}{\left(\cosh v-\cos u\right)^2}\nonumber\\
&=& \frac{2a^2 (\cosh ^{2}v-\cosh v\cos u)}{\cosh v-\cos u} \\ &=& \frac{2a^2 \cosh v}{\cosh v-\cos u} \label{toroidales2}
\end{eqnarray}
Con las ecuaciones (\ref{toroidales1}) y (\ref{toroidales2}) y se establece la siguiente relación
\begin{equation}
x^2+y^2+z^2+a^2=2a\coth v\sqrt{x^2+y^2}
\end{equation}
Esta expresión tiene la forma de la ecuación \ref{toroidales3},
que representa un toro.
\end{proof}

\begin{prop}
En coordenadas toroidales con valores constantes de $u$ se determina una esfera.
\end{prop}

\begin{proof}
Con las siguientes relaciones
{\small
\begin{eqnarray*}
x^{2}+y^{2} &=&\left( \frac{a\sinh v\cos \phi }{\cosh v-\cos u}\right)
^{2}+\left( \frac{a\sinh v\sin \phi }{\cosh v-\cos u}\right)
^{2} \\ &=& \frac{a^{2}\sinh ^{2}v}{\left( \cos u-\cosh v\right) ^{2}} \\
\left( z-a\cot u\right) ^{2} &=& \left( \frac{a\sin u}{\cosh v-\cos u}-\frac{a\cos u}{\sin u}\right) ^{2} % =\frac{a^{2}\left( \sin ^{2}u-\left( \cosh v-\cos u\right) \cos u\right) ^2}{\left( \left( \cosh v-\cos u\right) \sin u\right) ^{2}}\\
= \frac{a^{2}\left( 1-\cos u\cosh v\right) ^{2}}{\left(
\left( \cosh v-\cos u\right) \sin u\right) ^{2}}
\end{eqnarray*}
}

Se sigue que
\begin{eqnarray}
x^{2}+y^{2}+\left( z-a\cot u\right) ^{2}
% &=& \frac{a^{2}\sinh ^{2}v}{\left( \cos u-\cosh v\right) ^{2}}+\frac{a^{2}\left( 1-\cos u\cosh v\right) ^{2}}{\left( \left( \cosh v-\cos u\right) \sin u\right) ^{2}} \nonumber \\
&=& \frac{a^{2}\left( \cos u-\cosh v\right) ^{2}}{\left( \left( \cosh v-\cos u\right)
\sin u\right) ^{2}} =\frac{a^{2}}{\sin ^{2}u} \label{toroidales4}
\end{eqnarray}
La ecuación \ref{toroidales4} representa un esfera para valores constantes de $u$.
\end{proof}

Estas superficies se obtienen al ejecutar las siguientes instrucciones

%

\begin{tcolorbox}[title=\bf Instrucción \verb"Mathematica"]\begin{Verbatim}[formatcom=\letraverbatim]
x=2*Sinh[v]*Cos[\[Phi]]/(Cosh[v]-Cos[u]);
y=2*Sinh[v]*Sin[\[Phi]]/(Cosh[v]-Cos[u]);
z=2*Sin[u]/(Cosh[v]-Cos[u]);
\[Phi]=Pi/3;
A1=ParametricPlot3D[{x,y,z},{u,0,2*Pi},{v,-1,1},
MeshStyle->Black,PlotStyle->Orange,Mesh->30];
u=Pi/6;
A2=ParametricPlot3D[{x,y,z},{\[Phi],0,2*Pi},{v,0,2*Pi},
MeshStyle->Gray,PlotPoints->50,PlotStyle->Cyan];
v=Pi/4;
A3=ParametricPlot3D[{x,y,z},{\[Phi],0,2*Pi},{u,0,2*Pi},
MeshStyle->Gray,PlotPoints->50];
Show[A1,A2,A3,PlotRange->{{-5,5},{-5,5},{-6,7}}]
\end{Verbatim}
\end{tcolorbox}

Cuyo producto es el siguiente.

\begin{figure}[H]
\centering
\caption{Superficies constantes en coordenadas toroidales}
\label{fig:toroidales_constantes}
\includegraphics[width=0.30\textwidth]{Img/T_57a.pdf}
\includegraphics[width=0.30\textwidth]{Img/T_57b.pdf}

\fuentepropia
\end{figure}

\subsubsection{Coordenadas esferoidales oblatas.} \index{Coordenadas!esferoidales oblatas}

Un esferoide oblato es el sólido que se obtiene al hacer una elipse sobre su eje menor.
Las coordenadas esferoides oblatas (achatadas) se originan cuando una elipse rota alrededor de su eje menor,
la relación con las coordenadas cartesianas es la siguiente.

\begin{equation}
x = a \cosh\eta\cos\theta\cos\phi,\hskip0.1cm y = a \cosh\eta\cos\theta\sin\phi, \hskip0.1cm z= a\sinh\eta\sin\theta
\end{equation}
%
Con $\eta \in \mathbb{R}$, $0\leq \theta <2\pi$ y $0\leq \phi <\pi$.
Con la definición se pueden demostrar las siguientes identidades:

{\small
\begin{eqnarray*}
\frac{x^{2}+y^{2}}{a^{2}\cosh ^{2}\eta }+\frac{z^{2}}{a^{2}\sinh ^{2}\eta} &=&  \frac{a^{2}\cos ^{2}\theta \cosh ^{2}\eta }{a^{2}\cosh ^{2}\eta }+\frac{a^{2}\sinh ^{2}\eta \sin ^{2}\theta}{a^{2}\sinh ^{2}\eta }= 1 \\[0.1cm]
\frac{x^{2}+y^{2}}{a^{2}\cos ^{2}\theta }-\frac{z^{2}}{a^{2}\sin ^{2}\theta} &=& \frac{a^{2}\cos ^{2}\theta \cosh ^{2}\eta }{a^{2}\cos ^{2}\theta }-\frac{a^{2}\sinh ^{2}\eta \sin ^{2}\theta }{a^{2}\sin^{2}\theta }= 1.
\end{eqnarray*}
}

Es posible ver que para valores constantes de $\eta$ la superficie generada es un elipsoide,
mientras que para valores constantes de $\theta$ se obtienen hiperboloides de un manto con $z$
como eje de simetría. De la identidad $\tan \phi = y/x$ se sigue que para valores constantes de $\phi$
las superficies obtenidas son semiplanos. Presentamos algunas de estas superficies.

\begin{figure}[H]
\centering
\caption{Superficies constantes en coordenadas esferoidales oblatas}

\includegraphics[width=3cm,height=3.9cm]{Img/T_58a.pdf}

\includegraphics[width=3cm,height=3.9cm]{Img/T_58b.pdf}
\fuentepropia
\end{figure}

Con las siguientes instrucciones se obtienen estas figuras:

\begin{tcolorbox}[title=\bf Instrucción \verb"Mathematica"]\begin{Verbatim}[formatcom=\letraverbatim]
x=Cosh[\[Eta]]*Cos[\[Theta]]*Cos[\[Phi]];
y=Cosh[\[Eta]]*Cos[\[Theta]]*Sin[\[Phi]];
z=Sinh[\[Eta]]*Sin[\[Theta]];
\[Eta]=3/2;
A1=ParametricPlot3D[{x,y,z},{\[Theta],0,2*Pi},
{\[Phi],0,Pi}];
\[Theta]=Pi/6;
A2=ParametricPlot3D[{x,y,z},{\[Eta],-2,2},
{\[Phi],0,2*Pi},
PlotStyle->Orange];
\[Phi]=Pi/4;
A3=ParametricPlot3D[{x,y,z},{\[Eta],-2,2},
{\[Theta],0,Pi},PlotStyle->Cyan];
Show[A3,A2,A1,BoxRatios->{1,1,1}]
\end{Verbatim}
\end{tcolorbox}

\subsubsection{Coordenadas esferoides prolatas} \index{Coordenadas!esferoides prolatas}

Al hacer girar una elipse sobre su eje mayor el sólido obtenido se denomina
esferoide prolato.
Las coordenadas esferoides prolatas se obtiene cuando una elipse rota alrededor de su eje mayor.
La relación de este sistema de coordenadas con las coordenadas cartesianas es la siguiente:

\begin{equation}
x = a\sinh \eta \sin \theta \cos \phi,\hskip0.1cm y = a\sinh \eta \sin \theta \sin \phi, \hskip0.1cm z=a\cosh \eta \cos \theta
\end{equation}
con $\eta \in \mathbb{R}$, $0\leq \theta <\pi$ y $0\leq \phi < 2\pi$.
De las siguientes igualdades se evidencia que para valores constantes de $\eta$ la superficie
generada es un elipsoide, mientras que para valores constantes de $\theta$ se obtienen hiperboloides
de dos mantos con $z$ como eje de simetría.
De la identidad $\tan \phi = y/x$ se sigue que para valores constantes de $\phi$ las superficies obtenidas son semiplanos.
{\small
\begin{eqnarray*}
\frac{x^{2}+y^{2}}{a^{2}\sinh ^{2}\eta}+\frac{z^{2}}{a^{2}\cosh^{2}\eta} &=&
\frac{a^{2}\sin ^{2}\theta \sinh ^{2}\eta }{a^{2}\sinh ^{2}\eta}+\frac{a^{2}\cosh^{2}\eta\cos^{2}\theta}{a^{2}\cosh^{2}\eta}= 1 \\[0.2cm]
\frac{z^{2}}{a^{2}\cos ^{2}\theta }-\frac{x^{2}+y^{2}}{a^{2}\sin ^{2}\theta}
&=& \frac{a^{2}\cosh ^{2}\eta \cos ^{2}\theta }{a^{2}\cos^{2}\theta }-\frac{a^{2}\sin ^{2}\theta \sinh ^{2}\eta }{a^{2}\sin^{2}\theta }= 1
\end{eqnarray*}
}

\begin{figure}[H]
\centering
\caption{Superficies constantes en coordenadas esferoidales prolatas}
\includegraphics[width=4cm,height=4cm]{Img/T_59a.pdf}
\includegraphics[width=4cm,height=4cm]{Img/T_59b.pdf}
\fuentepropia
\end{figure}

Con las siguientes instrucciones se obtienen estas figuras:

\begin{tcolorbox}[title=\bf Instrucción \verb"Mathematica"]\begin{Verbatim}[formatcom=\letraverbatim]
x=Sinh[\[Eta]]*Sin[\[Theta]]*Cos[\[Phi]];
y=Sinh[\[Eta]]*Sin[\[Theta]]*Sin[\[Phi]];
z=Cosh[\[Eta]]*Cos[\[Theta]];
\[Eta]=3/2;
A1=ParametricPlot3D[{x,y,z},{\[Theta],0,Pi},{\[Phi],
0,2*Pi},MeshStyle->Black,PlotStyle->Orange,Mesh->30];
\[Theta]=Pi/6;
A2=ParametricPlot3D[{x,y,z},{\[Eta],-2,2},{\[Phi],0,
2*Pi},MeshStyle->Gray,PlotPoints->50,PlotStyle->Cyan];
\[Phi]=Pi/4;
A3=ParametricPlot3D[{x,y,z},{\[Eta],-2,2},{\[Theta],
0,Pi},PlotStyle->Cyan,PlotPoints->50];
Show[A1,A3]
\end{Verbatim}
\end{tcolorbox}

\subsubsection{Coordenadas cónicas} \index{Coordenadas!cónicas}

Las siguientes ecuaciones definen este sistema de coordenadas:
{\small
\begin{equation*}
x = \frac{uvw}{ab},\hskip0.1cm y = \frac{u}{a}\sqrt{\frac{\left( v^2-a^2\right) \left(
w^2-a^2\right) }{a^2-b^2}}, \hskip0.1cm z=\frac{u}{b}\sqrt{\frac{\left( v^2-b^2\right)
\left(w^2-b^2\right) }{b^2-a^2}}
\end{equation*}
}

Se debe considerar la siguiente condición $b^2>v^2>a^2>w^2$.
En este sistema de coordenadas, valores constante de $u$ generan esferas;
esto es consecuencia de la siguiente igualdad.
{\small
\[ x^2+y^2+z^2 = \frac{u^2v^2w^2}{a^2b^2}+\frac{u^2}{a^2}\frac{%
\left( a^2-v^2\right) \left( a^2-w^2\right) }{a^2-b^2}%
- \frac{u^2}{b^2}\frac{\left( b^2-v^2\right) \left(
b^2-w^2\right) }{\left( a^2-b^2\right) }= u^2\]
}

Presentamos ahora otras dos implicaciones.

\begin{prop}
En coordenadas cónicas con valores constante de $v$ se representan conos el\'{\i}pticos
con eje de simetr\'{\i}a el eje $z$.
\end{prop}

\begin{proof}
\begin{multline*}
\frac{x^2}{v^2}+\frac{y^2}{v^2-a^2}+\frac{z^2}{v^2-b^2} \\
= \frac{u^2w^2}{a^2b^2}-\frac{1}{a^2}\frac{u^2\left( a^2-w^2\right) }{a^2-b^2}+\frac{1}{b^2}\frac{u^2\left( b^2-w^2\right) }{a^2-b^2} =0.
\end{multline*}
Por lo anteriormente expuesto y considerando que el término $ v^2-b^2$ es negativo,
se concluye que  $\frac{x^2}{v^2}+\frac{y^2}{v^2-a^2}=\frac{z^2}{b^2-v^2}$.
Esta ecuación representa un cono elíptico, con eje simetría el eje $z$.
\end{proof}

\begin{prop}
En coordenadas cónicas, para valores constantes de $w$ la superficie
representada es un cono elíptico con el eje $x$
como eje de simetr\'{\i}a.
\end{prop}

\begin{proof}
\begin{multline*}
\frac{x^2}{w^2}+\frac{y^2}{w^2-a^2}+\frac{z^2}{w^2-b^2}
= \frac{u^2v^2}{a^2b^2}-\frac{1}{a^2}\frac{u^2\left( a^2-v^2\right)}{a^2-b^2}
+\frac{1}{b^2}\frac{u^2\left(b^2-v^2\right) }{a^2-b^2} = 0
\end{multline*}
Considerendo que los términos $w^2-a^2$ y $w^2-b^2$ son negativos, la expresión
$\frac{y^2}{a^2-w^2}+\frac{z^2}{b^2-w^2}=\frac{x^2}{w^2}$,
describe un cono elíptico con $x$ como eje de simetr\'{\i}a y $w$ constante.

\end{proof}
\begin{ejer}
Obtener algunas superficies constantes en coordenadas cónicas.
\end{ejer}

\subsubsection{Coordenadas biesféricas} \index{Coordenadas!biesféricas}

Las expresiones que definen este sistema de coordenadas son
\begin{equation}
x = \frac{a\sin \theta \cos \phi }{\cosh \eta -\cos \theta} ,\hskip0.25cm y = \frac{a\sin \theta \sin \phi }{\cosh \eta -\cos \theta }, \hskip0.25cm z=\frac{a\sinh \eta }{\cosh \eta -\cos \theta}
\end{equation}
En donde $\eta \in \mathbb{R}$, $\theta \in [0,\pi] $, $\phi \in [0,2\pi] $.

\begin{prop}
En coordenadas biesféricas con valores constantes de $\eta $ se describen esferas.
\end{prop}

\begin{proof}
A partir de la expresión $x^2+y^2+\left( z-a\coth \eta \right) ^2$, se sigue
\begin{multline*}
\hskip-0.3cm \left( \frac{a\sin \theta\cos \phi }{\cosh \eta -\cos \theta }\right)^2+\left(\frac{a\sin \theta
\sin \phi }{\cosh \eta -\cos \theta }\right) ^2+\left( \frac{a\sinh \eta}{\cosh \eta -\cos \theta }-\frac{a\cosh \eta }{\sinh \eta }\right) ^2 \\
% =\frac{a^2}{\left( \cosh \eta -\cos \theta \right) ^2}\left( \sin^2\theta +\left( \frac{\sinh ^2\eta -\left( \cosh \eta-\cos\theta\right) \cosh \eta }{\sinh \eta }\right) ^2\right)  \\
= \frac{a^2}{\left( \cosh \eta -\cos \theta \right) ^2} \,\frac{\left(
\cosh \eta -\cos \theta \right) ^2}{\sinh ^2\eta } = \frac{a^2}{\sinh ^2\eta }
\end{multline*}
Por lo tanto, la ecuación $x^2+y^2+\left( z-a\coth \eta \right)^2= a^2 {\rm csch} ^2\eta$,
para valores constantes de $\eta $, representa una esfera.
\end{proof}

\begin{prop}
En coordenadas biesféricas se satisface la siguiente igualdad
\begin{equation}\label{biesfericas3}
x^2+y^2+z^2= a^2 +2a\sqrt{x^2+y^2}\cot\theta.
\end{equation}
\end{prop}

\begin{proof}
A partir de las siguientes igualdades
\begin{eqnarray}
\hskip-0.3cm x^2+y^2+z^2 % &=&\left(\frac{a\sin\theta\cos\phi}{\cosh\eta-\cos\theta}\right)^2+\left( \frac{a\sin \theta\sin\phi}{\cosh\eta-\cos\theta}\right)^2+\left(\frac{a\sinh\eta}{\cosh\eta-\cos\theta}\right)^2 \nonumber \\
&=& \frac{a^2(\left(\sin\theta \cos \phi\right)^2+\left( \sin \theta \sin \phi \right) ^2+\left( \sinh \eta \right) ^2)}{\left( \cosh \eta -\cos \theta \right)^2}\nonumber \\
& =& a^2\frac{\cosh ^2\eta -\cos ^2\theta }{\left( \cosh \eta -\cos \theta \right) ^2} = a^2\frac{\left( \cosh \eta +\cos \theta\right) }{\left( \cosh \eta -\cos \theta \right) }  \label{biesfericas1}
\end{eqnarray}
Además
\begin{eqnarray}
\sqrt{x^2+y^2} &=& \sqrt{\left( \frac{a\sin \theta \cos \phi }{\cosh \eta-\cos \theta }\right) ^2+\left( \frac{a\sin \theta \sin \phi }{\cosh \eta -\cos \theta }\right) ^2} \nonumber  \\
&=& \sqrt{\frac{a^2\sin ^2\theta }{\left( \cosh \eta -\cos \theta \right) ^2}} =\frac{a\sin \theta }{\cosh \eta -\cos \theta } \label{biesfericas2}
\end{eqnarray}
Con las expresiones \ref{biesfericas1} y \ref{biesfericas2}
se establece que  % $x^2+y^2+z^2-2a\sqrt{x^2+y^2}\cot\theta= a^2$. por lo tanto
{\small
\begin{eqnarray*}
x^2+y^2+z^2-2a\sqrt{x^2+y^2}\cot \theta &=& a^2\frac{\left(\cosh \eta +\cos \theta \right) }{\left(\cosh \eta -\cos \theta \right) } -\frac{2a^2\sin \theta\cot \theta}{\cosh \eta -\cos \theta } \\
& =& a^2
\end{eqnarray*}
}
\end{proof}
A partir de relación \ref{biesfericas3} se observa que para $\theta <\pi/2$,
la superficie generada tiene forma de \emph{manzana} y
para $\theta >\pi/2$ se genera una superficie posee forma de \emph{limón.}

\begin{figure}[H]
\centering
\caption{Superficies constantes en coordenadas esferoidales prolatas}
\includegraphics[width=4.2cm,height=3.8cm]{Img/T_60a.pdf}
\includegraphics[width=4.2cm,height=3.8cm]{Img/T_60b.pdf}
\fuentepropia
\end{figure}
Con las siguientes instrucciones se obtienen estas figuras:

\begin{tcolorbox}[title=\bf Instrucción \verb"Mathematica"]\begin{Verbatim}[formatcom=\letraverbatim]
w=Cosh[\[Eta]]-Cos[\[Theta]];
{x,y,z}={Sin[\[Theta]]*Cos[\[Phi]]/w,
Sin[\[Theta]]*Sin[\[Phi]]/w,Sin[\[Eta]]/w}
\[Phi]=Pi/3;
A1=ParametricPlot3D[{x,y,z},{\[Theta],0,2*Pi},
{\[Eta],-2,2},PlotStyle->Cyan,Mesh->60];
\[Eta]=1/2;
A2=ParametricPlot3D[{x,y,z},{\[Theta],0,Pi},
{\[Phi],0,2*Pi},PlotPoints->50,PlotStyle->Orange];
\[Theta]=3*Pi/2;
A3=ParametricPlot3D[{x,y,z},{\[Phi],0,2*Pi},
{\[Eta],-1,1},PlotStyle->Cyan,PlotPoints->50];
Show[A3,A1,A2,PlotRange->{-2,4}]
\end{Verbatim}
\end{tcolorbox}

\subsubsection{Coordenadas cardioidicas} \index{Coordenadas!cardioidicas}

En este sistema de coordenadas se generan, en algunos, casos cardioides de revolución.

\begin{figure}[H]
\centering
\caption{Superficies constantes en coordenadas esferoidales prolatas}
\includegraphics[height=6cm]{Img/T_61a.pdf}\hskip0.5cm
\includegraphics[height=6cm]{Img/T_61b.pdf}
\fuentepropia
\end{figure}

Las expresiones que lo determinan son
\begin{equation}
x =\frac{uv\cos \theta}{\left( u^2+v^2\right)^2} ,\hskip0.5cm y =\frac{uv\sin \theta}{\left( u^2+v^2\right)^2}, \hskip0.5cm z=\frac{1}{2}\frac{u^2-v^2}{\left(u^2+v^2\right)^2}
\end{equation}
Con $u\in [0,\infty )$, $v\in [0,\infty)$, $\theta \in [0,2\pi)$.
Para valores de constante de $u$ y $v$ se generan
cardioides de revolución.

Estas figuras se obtienen utilizando las siguientes instrucciones:

\begin{tcolorbox}[title=\bf Instrucción \verb"Mathematica"]\begin{Verbatim}[formatcom=\letraverbatim]
x=u*v*Cos[\[Theta]]/(u^2+v^2)^2;
y=u*v*Sin[\[Theta]]/(u^2+v^2)^2;
z=(u^2-v^2)/(2*(u^2+v^2)^2);
u=1/2;
A1=ParametricPlot3D[{x,y,z},{\[Theta],0,2*Pi},
{v,0,4},PlotPoints->50];
v=1/3;
A2=ParametricPlot3D[{x,y,z},{\[Theta],0,2*Pi},{u,0,2},
MeshStyle->Gray,PlotStyle->RGBColor[0.2,0.6,0.4]];
\[Theta]=2*Pi/3;
A3=ParametricPlot3D[{x,y,z},{u,0,3},{v,0,3},
PlotStyle->Gray,PlotPoints->50];
Show[A2,A3,A1,PlotRange->{{-3,3},{-3,3},{-5,2}}]
\end{Verbatim}
\end{tcolorbox}

La ecuación $r=1+\cos\theta$ representa una cardioide.
A partir de ella se obtiene la siguiente expresión
$r=1+\cos\theta \to r^2=r+r\cos\theta \to x^2+y^2=\sqrt{x^2+y^2}+x;$ %$
usaremos esta misma relación con tres variables.

\begin{prop}
En coordenadas cardioidicas, para valores constantes de $u$ y $v$
se generan cardioides de revolución alrededor del eje $z$ y para valores constantes de
$\theta$, las superficies generadas son semiplanos.
\end{prop}
%
\begin{proof}
A partir de la siguiente relación
\begin{eqnarray*}
x^2+y^2+z^2 &=& \left( \frac{uv\cos \theta}{\left( u^2+v^2\right) ^2} \right) ^2+\left( \frac{uv\sin \theta}{\left( u^2+v^2\right) ^2} \right) ^2+\left( \frac{1}{2}\frac{u^2-v^2}{\left(
u^2+v^2\right) ^2}\right) ^2 \\ &=& \frac{1}{4\left(u^2+v^2\right) ^2}
\end{eqnarray*}

Se pueden obtener las siguientes ecuaciones
\begin{eqnarray*}
\sqrt{x^2+y^2+z^2}+z &=& \frac{1}{2\left( u^2+v^2\right) }+\frac{1}{2}%
\frac{u^2-v^2}{\left( u^2+v^2\right) ^2}= \frac{u^2}{%
\left( u^2+v^2\right) ^2} \\
\sqrt{x^2+y^2+z^2}-z &=& \frac{1}{2\left( u^2+v^2\right) }-\frac{1}{2}%
\frac{u^2-v^2}{\left( u^2+v^2\right) ^2}= \frac{v^2}{%
\left( u^2+v^2\right) ^2}
\end{eqnarray*}

Estas expresiones permiten establecer las siguientes igualdades
\[x^2+y^2+z^2=\frac{1}{4u^2}\left( \sqrt{x^2+y^2+z^2}+z\right)\]
\[x^2+y^2+z^2=\frac{1}{4v^2}\left( \sqrt{x^2+y^2+z^2}-z\right)\]

Estas ecuaciones corresponden a cardioides de revolución.
De la expresión $\tan \theta =\frac{y}{x}$ se observa que valores constantes $\theta$
determinan semiplanos.
\end{proof}

\subsubsection{Coordenadas elipsoidales.} \index{Coordenadas!elipsoidales}
Las coordenadas elipsoidales confocales $\left( u,v,w\right) $ o simplemente elipsoidales se relacionan con el sistema de coordenadas cartesianas como sigue
\begin{eqnarray*}
x^{2} &=&\frac{u^2 v^2 w^2}{b^2 c^2} \\
y^{2} &=&\frac{(u^2-b^2)(v^2-b^2)(b^2-w^2)}{b^2(c^2-b^2)} \\
z^{2}&=& \frac{(u^2-c^2)(c^2-v^2)(c^2-w^2)}{c^2(c^2-b^2)}
\end{eqnarray*}
En donde  $0<w^2<b^2<v^2<c^2<u^2$. Este sistema de coordenadas satisface las siguientes identidades.
\begin{eqnarray}
\frac{x^2}{u^2}+\frac{y^2}{u^2-b^2}+\frac{z^2}{u^2-c^2}&=& 1 \label{elipsoidales1}\\
\frac{x^2}{v^2}+\frac{y^2}{v^2-b^2}+\frac{z^2}{v^2-c^2} &=&1 \label{elipsoidales2}\\
\frac{x^2}{w^2}+\frac{y^2}{w^2-b^2}+\frac{z^2}{w^2-c^2} &=&1 \label{elipsoidales3}
\end{eqnarray}
De la ecuación \ref{elipsoidales1} se observa que si $u$ es constante la superficie obtenida es un elipsoide.
Con $v$ constante, la ecuación \ref{elipsoidales2} corresponde a un hiperboloide de un manto y si
$w$ es constante entonces la ecuación \ref{elipsoidales3} determina un hiperboloide de dos mantos.
\begin{ejer}
Hallar algunas superficies constantes en coordenadas elipsoidales.
\end{ejer}

\subsubsection{Coordenadas esféricas y cil\'{\i}ndricas} \index{Coordenadas!esféricas} \index{Coordenadas!cilíndricas}
Estas coordenadas son quizás las más estudiadas en un curso de cálculo vectorial. En la página \pageref{esfericas} las consideramos, recordemos su definición
\[x=\rho \cos \theta \sin \phi, \, y=\rho \sin \theta \sin \phi,\,z=\rho \cos \phi\]
con $\rho \geq 0$, $\theta \in \lbrack 0,2\pi ]$, $\phi \in \lbrack 0,\pi ]$,
De la relación
\[x^{2}+y^{2}+z^{2}=(\rho \cos \theta \sin \phi)^{2}+(\rho \sin \theta \sin \phi)^{2}+(\rho \cos \phi) ^{2}=\rho ^{2}\]
se infiere que para valores constantes de $\rho $ las superficies obtenidas son esferas.
Al simplificar
\[x\sin \theta =( \rho \cos \theta \sin \phi)\sin \theta =( \rho \sin \theta \sin \phi ) \cos \theta
=y\cos \theta\]
se evidencia que si $\theta $ toma valores constantes se obtienen semiplanos perpendiculares la plano polar;
la equivalencia
\[(x^{2}+y^{2}) \cos ^{2}\phi = ( ( \rho \cos \theta\sin \phi ) ^{2}+( \rho \sin \theta \sin \phi ) ^{2})
\cos^{2}\phi =z^{2}\sin^{2}\phi\]
permiten concluir que si $\phi $ es constante se expresiones que determinan semiconos.
Las figuras de algunas superficies se obtienen ejecutando las siguientes instrucciones:

\begin{tcolorbox}[title=\bf Instrucción \verb"Mathematica"]\begin{Verbatim}[formatcom=\letraverbatim]
x=\[Rho]*Cos[\[Theta]]*Sin[\[Phi]];
y=\[Rho]*Sin[\[Theta]]*Sin[\[Phi]];
z=\[Rho]*Cos[\[Phi]];
\[Phi]=Pi/3;
A1=ParametricPlot3D[{x,y,z},{\[Theta],0,2*Pi},
{\[Rho],0,3}];
\[Theta]=Pi/6;
A2=ParametricPlot3D[{x,y,z},{\[Phi],0,Pi},
{\[Rho],0,3}];
\[Rho]=2;
A3=ParametricPlot3D[{x,y,z},{\[Phi],0,Pi},
{\[Theta],0,2*Pi},PlotStyle->Cyan];
Show[A3,A2,A1,PlotRange->{{-3,3},{-3,3},{-3,3}}]
\end{Verbatim}
\end{tcolorbox}

El producto obtenido es el siguiente.

\begin{figure}[H]
\centering
\caption{Superficies constantes en coordenadas esféricas}
\includegraphics[height=3.6cm]{Img/T_62a.pdf}
\includegraphics[height=3.6cm]{Img/T_62b.pdf}
\includegraphics[height=3.6cm]{Img/T_62c.pdf}
\fuentepropia
\end{figure}

En coordenadas cil\'{\i}ndricas se establece que
$ x=r\cos \theta, \, y=r\sin \theta,\, z=z$,
con $r\geq 0$, $\theta \in \lbrack 0,2\pi ]$ y $z \in \mathbb{R}$.
De las relaciones $x^{2}+y^{2}=r^{2}$ se infiere que para valores constantes
de $r$ las superficies obtenidas son cilindros; Si
$ x\sin \theta =r\cos \theta \sin \theta =( r\sin \theta ) \cos \theta =y\cos \theta $, se obtiene ecuaciones
que representan semiplanos paralelos al eje $z$. Las ecuaciones
$z=k$ representan planos paralelos al plano polar.
